%% CPP 2027 Submission — 0pat Exercise XXXIII: The Imaginary Part Compressed into a Coefficient Axis
\documentclass[sigplan,10pt,anonymous,review]{acmart}
\settopmatter{printfolios=true,printacmref=false}

\AtBeginDocument{%
  \providecommand\BibTeX{{\normalfont B\scshape ib\TeX}}}

\settopmatter{printacmref=false}
\renewcommand\footnotetextcopyrightpermission[1]{}
\pagestyle{plain}

\usepackage{microtype}
\usepackage{booktabs}
\usepackage{amsmath}
% XeTeX (tectonic): acmart loads newtxmath/libertine which already defines \Bbbk;
% reset it so amssymb's definition is accepted.
\let\Bbbk\relax
\usepackage{amssymb}

\begin{document}

\title{The Imaginary Part Compressed into a Coefficient Axis: Where the\\
Product of Imaginary Parts Appears (Exercise XXXIII)}

\author{Anonymous}
\affiliation{%
  \institution{Anonymous}
}
\email{anonymous@example.org}

\begin{abstract}
We formalize what happens when the imaginary part is compressed into a
coefficient axis. In the complex plane, the imaginary axis does not
represent the imaginary part itself; it represents the real
coefficients of the imaginary part. The unit $i$ is pressed into the
axis's semantics. We prove, with five machine-checked theorems, the
consequence of this compression: in multiplication, the product of two
imaginary-part coefficients $b d$ leaves the imaginary part and appears
in the real part, with a sign determined by the square of the unit. In
$\mathbb{C}$ ($i^2 = -1$): the real part of $(a+bi)(c+di)$ is
$ac - bd$. In the split-complex frame ($j^2 = +1$): the real part of
$(a+bj)(c+dj)$ is $ac + bd$. The cross terms $ad + bc$ stay in the
imaginary part in both. The leak sign is the unit square: $i \cdot i =
-1$ versus $j \cdot j = +1$. Zero errors, zero \texttt{sorry}.
\end{abstract}

\keywords{Lean, formalization, complex multiplication, coefficient compression, product leak, 0pat}

\maketitle

\section{The compression}

The imaginary axis of the complex plane displays, at each tick, a real
number $t$ --- the coefficient of the imaginary part. The imaginary
part itself, $t \cdot i$, is the product of this visible coefficient and
the invisible unit $i$. The unit is compressed into the axis: the axis
shows coefficients, not units.

This exercise proves what happens under this compression, in
multiplication. Two numbers $a + bi$ and $c + di$ multiply to
$(ac - bd) + (ad + bc)i$. The observation: the product of the two
imaginary-part coefficients, $b d$, does \emph{not} stay in the
imaginary part. It appears in the real part, with a sign. That sign is
the square of the unit.

\section{The theorems}

\paragraph{The axis shows coefficients (complex\_imag\_axis\_coefficient).}
A point $z$ of the imaginary axis satisfies $z = t \cdot i$ exactly
when $z.\mathrm{re} = 0$ and $z.\mathrm{im} = t$. The tick $t$ is a
real coefficient; the unit $i$ is not displayed by the axis.

\paragraph{The product of imaginary parts leaks to the real part
(complex\_mul\_imag\_leaks).} The real part of $(a+bi)(c+di)$ is
$ac - bd$. The term $bd$ is the product of the two imaginary-part
coefficients; it appears in the real part with a minus sign, because
$i \cdot i = -1$.

\paragraph{The cross terms stay in the imaginary part
(complex\_mul\_cross\_stays\_imag).} The imaginary part is
$ad + bc$: one imaginary coefficient times one real coefficient stays
in the imaginary part. Only the product of two imaginary coefficients
leaks.

\paragraph{The split frame leaks with the opposite sign
(split\_mul\_imag\_leaks).} In the split-complex frame ($j^2 = +1$),
the real part of $(a+bj)(c+dj)$ is $ac + bd$. The same compression, the
same leak --- but the sign is now plus, because $j \cdot j = +1$.

\paragraph{The leak sign is the unit square
(leak\_sign\_is\_unit\_square).} $i \cdot i = -1$ in $\mathbb{C}$ and
$j \cdot j = +1$ in the split frame. The sign with which the product of
imaginary coefficients appears in the real part is exactly the square
of the unit that was compressed into the axis.

\section{Honest boundary and position}

The content is elementary algebra (multiplication formulas); the value
is the precise statement of the compression's consequence. The result
concerns the coordinate structure of the frames; no claim is made about
the Riemann zeta function. The analysis of $\zeta$ remains the open
span of the bridge.

\section{Formalization}

The proof is a single Lean file (\texttt{proof.lean}, 95 lines),
importing \texttt{Mathlib.Data.Complex.Basic} and
\texttt{Mathlib.Data.Real.Basic}. Five theorems, compiled with Lean
4.32.2 / mathlib v4.32.2, zero errors, zero \texttt{sorry}.

\begin{center}
\begin{tabular}{lll}
\toprule
Theorem & Statement & Proof core \\
\midrule
\texttt{complex\_imag\_axis\_coefficient} & $z = t \cdot i \leftrightarrow \mathrm{re}\, z = 0 \wedge \mathrm{im}\, z = t$ & Complex.ext; simp \\
\texttt{complex\_mul\_imag\_leaks} & $\mathrm{re}((a+bi)(c+di)) = ac - bd$ & simp \\
\texttt{complex\_mul\_cross\_stays\_imag} & $\mathrm{im}((a+bi)(c+di)) = ad + bc$ & simp \\
\texttt{split\_mul\_imag\_leaks} & $\mathrm{re}((a+bj)(c+dj)) = ac + bd$ & rfl \\
\texttt{leak\_sign\_is\_unit\_square} & $i \cdot i = -1$, $j \cdot j = +1$ & ext; norm\_num \\
\bottomrule
\end{tabular}
\end{center}

\section{Conclusion}

The compression of the imaginary part into a coefficient axis has one
precise, machine-checked consequence: in multiplication, the product of
two imaginary-part coefficients leaks into the real part, and the sign
of the leak is the square of the unit compressed into the axis ($-1$
for $\mathbb{C}$, $+1$ for the split frame). The cross terms stay in
the imaginary part. The frames are fully understood; the analysis of
$\zeta$ remains the open span of the bridge.

\bibliographystyle{ACM-Reference-Format}
\bibliography{refs}

\end{document}
